\documentclass[11pt]{article}
\usepackage[margin=2.3cm]{geometry}
\usepackage[T1]{fontenc}
\usepackage[utf8]{inputenc}
\usepackage{lmodern}
\usepackage{enumitem}
\usepackage{booktabs}
\usepackage[colorlinks=true,linkcolor=blue!50!black,urlcolor=blue!50!black]{hyperref}
\usepackage{xcolor}
\setlist{nosep,leftmargin=1.4em}
\setlength{\parskip}{4pt}\setlength{\parindent}{0pt}

\begin{document}

{\large\bfseries Revision summary: a three-module recipe makes SODE-Guard
accurate on rare attacks, robust under adaptive attack and certified by its own
randomness, validated over five seeds with a protocol built for security use.}

\smallskip
\textit{General comment to the reviewers and the Area Chair (Submission 13607)}

We thank the reviewers and the Area Chair. Their comments shaped this revision,
and the uploaded manuscript reflects all of the work carried out during the
discussion.

\section*{What is new in the method}
\begin{enumerate}
\item \textbf{Imbalance-aware training}: a logit-adjusted loss, a full training
budget with early stopping on validation macro-F1, and a post-hoc logit
adjustment tuned on a held-out selection split.
\item \textbf{EOT-matched adversarial training}: the training adversary attacks
exactly the deployed eight-path predictor, with fresh Brownian paths at every
step (N-FGSM, $\varepsilon$ ramp, early stopping on robust validation macro-F1).
\item \textbf{A smoothing certificate that uses the model's own noise}.
Proposition~1 proves that randomised smoothing remains sound when the base
classifier is itself random. Gaussian input noise is a pure-diffusion prefix of
the SDE, so SODE-Guard is certified by sampling input noise and Brownian paths
jointly. Theorem~A remains as a deterministic Lipschitz certificate for the
discretised SDE.
\end{enumerate}

\section*{How it is evaluated}
\begin{itemize}
\item \textbf{Data:} container flows of ICS3D, with duplicates and conflicting
labels removed before splitting.
\item \textbf{Seeds and splits:} five seeds, each with a new 60/20/20 split.
\item \textbf{Attacks:} eight adaptive attacks (PGD with restarts,
APGD-CE/DLR with EOT, EOT-PGD, an expected-margin attack, Square, transfer),
each scored with fixed and with fresh Brownian paths.
\item \textbf{Checks:} gradient-masking checks, a full-strength re-attack, and a
per-class chronological split.
\item \textbf{Baselines:} LightGBM, XGBoost and a random forest.
\item \textbf{Statistics:} paired-bootstrap contrasts with Holm correction,
declared before the runs.
\end{itemize}

\section*{Main results (mean over five seeds)}
\begin{center}\small
\begin{tabular}{lccc}
\toprule
Metric & Before & With the recipe & Contrast \\
\midrule
Clean macro-F1 & 0.438 & 0.495 & $+0.057$, $p_{\mathrm{Holm}}=0.016$\\
Worst-case macro-F1 (eight attacks) & 0.245 & 0.420 & $+0.174$, $p_{\mathrm{Holm}}=0.016$\\
Certified accuracy, $\ell_2$ radius $0.273$ & 0.096 (Theorem~A) & 0.576 (0.877 test mix) & on par with smoothed MLP\\
\bottomrule
\end{tabular}
\end{center}
\begin{itemize}
\item Fixed and fresh Brownian paths agree (0.420 vs.\ 0.412), and certified
accuracy at radius 0.5 is 0.491.
\item SODE-Guard is on par with the strongest adversarially trained and smoothed
baselines (MLP and neural ODE trained with the same recipe).
\item Gradient-boosted trees reach the best clean macro-F1 (0.62) but fall to
0.02--0.04 under a black-box attack, which shows where robust neural detectors
are needed.
\item The full-strength re-attack and the chronological split confirm every
conclusion.
\end{itemize}

\section*{How the revision answers each review}
\textbf{Reviewer vjJK.}
\begin{itemize}
\item Adaptive attacks are central to the evaluation, and the masking checks pass.
\item The released pipeline attacks and certifies exactly the predictor that is
deployed.
\item The anti-concentration regulariser is replaced by three modules, each with
a significant measured effect.
\end{itemize}
\textbf{Reviewer FiAv.}
\begin{itemize}
\item The certificate is applied to trained detectors and is competitive with
randomised smoothing of an MLP.
\item The stochastic model's distinctive strength is now explicit: its own paths
are part of the smoothing distribution.
\item The protocol follows Arp et al.\ and TESSERACT: deduplication, multiple
seeds, chronological evaluation, strong tree baselines.
\item The paper cites and discusses Liu et al.\ (CVPR 2020).
\end{itemize}
\textbf{Reviewer kdck.}
\begin{itemize}
\item Every attack is scored with fixed and with fresh paths.
\item Recent robust and certified baselines are included, and BARS and MARS are
discussed.
\item Figures~1--2 are redrawn around the method and the accuracy / robustness /
certification trade-off.
\item The motivation is now grounded in the measured brittleness of the most
accurate clean models.
\end{itemize}

\section*{Ablations}
Every module was ablated over five seeds:
\begin{itemize}
\item the post-hoc adjustment and the training budget account for the clean gain;
\item single-step N-FGSM matches five-step PGD at lower cost;
\item the consistency weight trades certified accuracy at small radii against
large radii.
\end{itemize}

\section*{Future work}
We will extend SODE-Guard along the lines the reviewers opened:
\begin{itemize}
\item flow-sequence context for the rarest attack classes;
\item feature-constrained and packet-level attacks;
\item anisotropic noise in the certificate, in the spirit of BARS and MARS;
\item a second container corpus with cross-dataset transfer;
\item conformal alarm thresholds on the held-out calibration split.
\end{itemize}

The anonymised repository contains the code, the resumable multi-seed driver, a
Colab notebook and every per-seed result file behind the tables and figures.

\medskip
\textbf{SODE-Guard and its open protocol give practitioners a ready path to
intrusion detectors that are robust by measurement and certified by
construction, and give researchers a reproducible benchmark for the next
generation of robust detectors.}

\end{document}
